%% ============================================================
\section{Local-Model Adjudication: What a 12\,GB GPU Can and Cannot Decide}
\label{sec:local-llm}

The residual false positives characterised in \cref{sec:limitations}
are, by construction, the cases SqC's static analyses cannot settle.
An obvious question is whether a language model running on commodity
local hardware can settle them---not as a mandatory filter, but as an
\emph{optional force-adder} for users who have a GPU but no
frontier-model access.  We measured this directly.

\subsection{Setup}

All measurements ran on a single NVIDIA RTX 3060 (12\,GB, CUDA
capability 8.6, 10.6\,GiB usable) via a userland \texttt{ollama}
install, with models at their default \texttt{Q4\_K\_M} quantisation.
Each finding was presented as the rule identifier, SqC's own diagnostic
message, and a 50-line source window (35 leading, 15 trailing) with the
flagged line marked; the model answered \texttt{TP} or \texttt{FP} with
one sentence of justification, at temperature 0 with a fixed seed.

The evaluation set is drawn from the adjudicated oracle described in
\cref{sec:methodology}, restricted to SQLite at the pinned commit.  One
property of that oracle governs the whole experiment: \textbf{94.6\% of
its 88,939 labels (84,116) were adjudicated by an LLM}, and only 4,823
are human-adjudicated.  Evaluating an LLM triage stage against
LLM-written labels measures agreement with its own labeller, so every
headline claim below is additionally reported on the human-labelled
subset alone.

\subsection{A global filter does not work}

We first evaluated triage as a global filter over a 200-finding sample
(60 TP, 140 FP; baseline precision 30.0\%) spanning roughly thirty
rules (\cref{tab:llm-global}).  Because the base rate is low, accuracy
is uninformative---a
reject-everything policy scores 70\%---so we report true-positive
retention and false-positive suppression separately, with Matthews
correlation (MCC) as the discrimination summary.

\begin{table*}[t]
\centering
\small
\caption{Triage as a global filter over a 200-finding sample (60 TP,
140 FP; baseline precision 30.0\%).  The three models fail in different
directions rather than succeeding to different degrees.}
\label{tab:llm-global}
\begin{tabular}{llrrrr}
\hline
Model & Kind & TP ret. & FP supp. & Prec. & MCC \\
\hline
Qwen2.5-Coder-7B  & code    & 100.0\% & 2.9\%  & 30.6\% & $+0.094$ \\
Qwen2.5-Coder-14B & code    & 20.0\%  & 98.6\% & 85.7\% & $+0.334$ \\
Qwen3-8B          & general & 1.7\%   & 96.4\% & 16.7\% & $-0.051$ \\
\hline
\end{tabular}
\end{table*}

The 7B code model assents to essentially every
claim; the general-purpose 8B model rejects essentially every claim,
and---because it discards true positives preferentially---leaves
precision \emph{below} the untriaged baseline.  Only the 14B code model
shows real signal, and it buys precision of 85.7\% by discarding 80\%
of the true positives.

Two controls matter.  First, at equal scale the code-specialised model
is decisively better than the general-purpose one ($+0.334$ versus
$-0.051$), so code specialisation is the right axis.  Second, the
apparent discrimination is largely a threshold effect: re-running the
14B with a burden-of-proof prompt that instructs it to reject only on
demonstrable grounds moved the operating point sharply (retention
20.0\%$\rightarrow$68.3\%, suppression 98.6\%$\rightarrow$37.9\%) while
MCC \emph{collapsed} from $+0.334$ to $+0.059$.  Prompt engineering
slides these models along a near-chance trade-off curve; it does not
add judgement.

\subsection{Per-rule, the picture inverts}

The pooled figures conceal the result that matters.  Screening the 14B
model per-rule---615 findings, class-balanced at 25 TP and 25 FP per
rule across the 13 rules with adequate labels---separates one rule
cleanly from the rest (\cref{tab:llm-per-rule}).

\begin{table*}[t]
\centering
\small
\caption{Per-rule discrimination, Qwen2.5-Coder-14B, class-balanced.
``Human-only'' repeats balanced accuracy on the human-adjudicated
subset.}
\label{tab:llm-per-rule}
\begin{tabular}{lrrrrl}
\hline
Rule & TP ret. & FP supp. & Bal.\ acc. & MCC & Human-only \\
\hline
MSC17-C & 96\% & 88\%  & 91.9\% & $+0.840$ & 92\% ($n{=}49$) \\
PRE11-C & 13\% & 100\% & 56.7\% & $+0.275$ & 50\% ($n{=}2$) \\
PRE01-C & 16\% & 92\%  & 54.0\% & $+0.123$ & 50\% ($n{=}2$) \\
EXP34-C & 5\%  & 96\%  & 50.6\% & $+0.030$ & --- \\
DCL00-C & 0\%  & 100\% & 50.0\% & $+0.000$ & 50\% ($n{=}9$) \\
DCL03-C & 4\%  & 96\%  & 50.0\% & $+0.000$ & 50\% ($n{=}5$) \\
MSC13-C & 0\%  & 100\% & 50.0\% & $+0.000$ & --- \\
DCL13-C & 0\%  & 96\%  & 48.0\% & $-0.143$ & 50\% ($n{=}4$) \\
MSC04-C & 12\% & 84\%  & 48.0\% & $-0.058$ & 61\% ($n{=}16$) \\
PRE12-C & 0\%  & 92\%  & 46.0\% & $-0.204$ & --- \\
PRE10-C & 20\% & 62\%  & 41.2\% & $-0.193$ & --- \\
MEM05-C & 4\%  & 76\%  & 40.0\% & $-0.288$ & 50\% ($n{=}11$) \\
PRE00-C & 8\%  & 72\%  & 40.0\% & $-0.260$ & 30\% ($n{=}7$) \\
\hline
\end{tabular}
\end{table*}

MSC17-C reaches 91.9\% balanced accuracy (MCC $+0.840$); every other
rule sits between 40\% and 57\%, and six have negative MCC.  MSC17-C's
result is not a labelling artefact---it reproduces at 92\% on the
human-adjudicated subset ($n{=}49$), the largest human-labelled cell in
the screen.

Part of the separation is explained by the rules' own definitions.
MSC17-C asks whether a \texttt{case}
section ends in \texttt{break}, \texttt{return}, \texttt{continue}, or
\texttt{goto}, or else carries a fallthrough comment.  That question is
closed over the syntactic neighbourhood of the finding: a 50-line
window contains all the evidence needed to answer it.  DCL13-C asks
whether a pointer parameter is never modified by the function and
should therefore be \texttt{const}---a question requiring whole-program
aliasing and call-graph reachability into every callee.  No 50-line
window can contain that evidence, and the model scores at chance.

This gives a necessary condition for enabling a local-model
adjudicator on a given rule---that its decision procedure be closed
over the finding's syntactic neighbourhood---which is a static property
of the rule rather than something that must be measured per rule.  As
the next subsection shows, it is not a sufficient one.

\subsection{Window locality is necessary but not sufficient}

The obvious reading of \cref{tab:llm-per-rule} is that the model fails
wherever the deciding evidence lies outside the window, and that
assembling a richer context would recover those rules.  That reading is
correct for the whole-program rules and \emph{wrong} for the
preprocessor family, and the distinction is worth stating precisely
because it bounds what better context can buy.

We checked the \texttt{PRE*} findings directly.  All 602 of them sit on
a \texttt{\#define} line: these are definition-side rules, so the macro
body that decides them is in the window by construction, and only 17
(2.8\%) have continuation lines running past the trailing edge of the
window.  The model therefore had the deciding evidence in front of it
and still scored 40--57\%.  PRE01-C in particular---``is this macro
parameter parenthesised in the replacement text?''---is as syntactically
local a question as MSC17-C's, and the model reaches only 54.0\% on it.

Window locality is thus a \emph{necessary} condition for local
adjudication, not a sufficient one, and the criterion we offered above
must be weakened accordingly: a rule whose decision needs whole-program
facts is reliably beyond these models, but a rule whose decision is
local is not thereby within them.

What separates MSC17-C from PRE01-C appears to be the character of the
false positives rather than the character of the rule.  MSC17-C's false
positives are mechanically stereotyped---the adjudication record
attributes most of them to the checker failing to see a \texttt{break}
nested inside a braced case body---so the model is recognising one
recurring analyser mistake rather than exercising judgement about the
rule.  We verified that a residue of this class survives in the current
build: of 126 MSC17-C findings in the latest SQLite run, 12 (10\%) are
cases whose braced body closes with \texttt{\};}, where the stray empty
statement becomes the section's last statement and masks the
\texttt{break} that precedes it.  We report this as a partial
explanation only; we have not established that a single mechanism
accounts for the rule's remaining false positives.

This suggests the more promising use of a local model is not as a
runtime adjudicator but as a development instrument---a cheap,
independent detector of stereotyped analyser bugs, run against the
existing oracle rather than against user code.  The context-assembly
direction remains open for the whole-program rules (DCL13-C, MSC04-C),
where SqC's inter-procedural summaries could supply callee facts no
window contains, but we have not yet tested it.

\subsection{Cost and disposition}

Throughput was 2.75\,s per finding single-stream.  At that rate the 298
MSC17-C findings in our latest real-world run adjudicate in roughly 14
minutes, and the 3,151 findings SqC itself flags
\texttt{requires\_manual\_review} in about 2.4 hours single-stream
(considerably less batched).  Applying the same pass to all 69,487
findings would take about 53 hours.  Local adjudication is therefore
viable as an opt-in, per-rule extension and not as a mandatory
whole-corpus stage.

We note one further caveat on SqC's own abstention signal: the findings
it marks \texttt{requires\_manual\_review} carry \emph{lower} measured
precision (5.6\%) than the findings it decides outright (18.6\%), so
that flag presently selects a predominantly-false-positive set rather
than a genuinely ambiguous one.  Its selection logic warrants review
before it is used to route work to any adjudicator, human or machine.

Finally, these results bound only \emph{locally-runnable} models.  The
published gains from LLM-assisted static-analysis
triage~\cite{zerofalse,pathfeasibility} come from agentic loops driving
frontier models; the one comparative study we are aware of reports that
a code-specialised model \emph{underperformed} general frontier models
in that setting~\cite{siftingnoise}, and no sub-15B numbers appear in
that literature.  Our contribution here is the sub-15B datapoint, and
the finding that within it the decisive variable is not model choice
but whether the rule's decision procedure fits in the window.

We also caution against the purpose-built C vulnerability classifiers
(VulBERTa, SecureFalcon) as an alternative backbone.  They are binary
classifiers over a code snippet: they emit vulnerable/not-vulnerable
and have no input channel for the analyser's specific claim, which is
precisely what triage must verify.  Their headline accuracies are
additionally drawn from synthetic or duplicate-heavy corpora; the
PrimeVul study found a state-of-the-art 7B model falling from 68.3\%
F1 on BigVul to 3.09\% F1 once the benchmark was
deduplicated and correctly labelled~\cite{primevul}.

\subsection{Threats to validity}

The screen covers one project (SQLite) at one commit, one quantisation,
and one context geometry.  Per-rule cells are 25 findings per class, so
individual rates carry roughly $\pm$10--20 percentage points of
sampling error; the MSC17-C/rest separation is far larger than that
interval, but the ordering \emph{among} the near-chance rules is not
meaningful.  Human-labelled cells are smaller still, and for four rules
absent entirely.  We report MSC17-C as a robust positive result and the
remaining twelve as a robust negative one; finer distinctions within
the negative group would require a larger and more balanced oracle.
